\subsection{Спектроскопическое различение конфигураций в ансамбле}
\label{sec:experiment-spectroscopy}

Полученные различия каналов позволяют предложить эксперимент по определению
преимущественного расположения КТ относительно МНЧ. При этом необходимо
разделить ориентацию осей стержней в лабораторной системе координат,
положение КТ на стержне и направление переходного диполя внутри КТ.
Линейный дихроизм характеризует прежде всего первую величину; сам по себе
он не показывает, присоединена ли КТ к вершине или боковой поверхности.
Предлагаемая ниже процедура объединяет несколько измерений. Она является
рекомендацией для проверки расчётных выводов; обратная задача для ансамбля
в настоящей работе численно не решалась.

\paragraph{Ориентация МНЧ.}
Первое измерение -- спектры пропускания или экстинкции при вращении
линейной поляризации относительно выделенного направления образца.
Для этого достаточно спектрофотометра с поляризатором и поворотным
держателем. Необходимо регистрировать продольную и поперечную плазмонные
полосы и измерять контрольную плёнку без частиц. Ориентацию золотых
стержней в растянутом поливиниловом спирте (ПВС) и её проявление
в поляризационных спектрах непосредственно исследовали van der Zande
и соавторы~\cite{vanderzande1999}, а ориентацию в объёме и на поверхности
полимера -- Murphy и Orendorff~\cite{murphy2005}.
Для плоского образца в линейном режиме угловую зависимость можно
представлять через постоянную составляющую и гармонику
$\cos 2(\theta-\theta_0)$, извлекая направление директора $\theta_0$
и оптическую анизотропию. Пересчёт последней в степень ориентационного
порядка требует учёта обеих поляризуемостей и распределения наклонов;
отношение высот двух полос само по себе не является долей ориентированных
частиц. Если доступен тёмнопольный микроскоп, отдельные стержни можно
проверять поляриметрией рассеяния: такой способ определения азимутальной
ориентации продемонстрирован в~\cite{simo2025}. Для частиц малого размера
его применимость определяется отношением сигнала к шуму.

\paragraph{Признаки положения КТ.}
После определения директора следует измерить КТ-селективный отклик при
поле вдоль и поперёк него. Доступный вариант -- поляризационно-разрешённые
спектры возбуждения фотолюминесценции (ФЛ): при изменении энергии
возбуждения регистрируется интегральная полоса излучения КТ с нормировкой
на падающий поток фотонов. Анализатор в канале регистрации сначала
убирают или суммируют две ортогональные компоненты, чтобы не смешивать
анизотропию возбуждения с анизотропией сбора излучения. Дополнительная
серия с анализатором затем характеризует поляризацию самой ФЛ.
Спектры гибрида сопоставляют с КТ без металла и МНЧ без КТ в той же среде.
Возможность разделять поляризационно-зависимое возбуждение и изменение
экситонного распада в собранных комплексах показана
в~\cite{cohenhoshen2012}; применение этого принципа к рассматриваемым
торцевым и боковым конфигурациям является предлагаемой здесь схемой.

Ожидаемые различия на рабочей частоте иллюстрирует
табл.~\ref{tab:spectral-signatures}. У торцевой конфигурации продольное
возбуждение существенно эффективнее поперечного. При боковом размещении
становятся важны азимут КТ вокруг стержня и различие радиального
и касательного поперечных полей. Поэтому для боковой группы требуется
угловая серия, а не только два положения поляризатора. При случайном
азимутальном расположении КТ соответствующие вклады усредняются.
Числа таблицы относятся к определённым переходным диполям из
раздела~\ref{sec:geometry}; без учёта ориентаций реальных переходов они
не являются готовой калибровкой ансамбля.

\begin{table}[tbp]
  \centering
  \caption{Расчётные признаки конфигураций при $g=\SI{1}{\nano\meter}$
  в FQS с многоосцилляторной дисперсией. $G_{\rm exc}$ определено
  на фиксированной частоте $E_L=\SI{2.042}{\electronvolt}$.
  Порог дан в $10^{-6}$~Дж/см$^2$; прочерк означает отсутствие
  разрешённого порога в выполненном сканировании.}
  \label{tab:spectral-signatures}
  \small
  \begin{tabular}{llrr}
    \toprule
    Положение КТ & Поле & $G_{\rm exc}$ & $\mathcal F_\eta$ \\
    \midrule
    Вершина & вдоль оси & 25.3 & 1.62 \\
    Вершина & поперёк оси & 0.066 & --- \\
    Бок & вдоль оси & 2.19 & 26.9 \\
    Бок & поперёк, радиально & 4.73 & 9.39 \\
    Бок & поперёк, касательно & 0.013 & --- \\
    \bottomrule
  \end{tabular}
\end{table}

Для сопоставления узких спектральных признаков нужен перестраиваемый
узкополосный источник: ширина рассчитанной линии составляет около
\SI{2.5}{\milli\electronvolt}, а сдвиг -- доли или единицы мэВ.
Предпочтительно разрешение существенно лучше ширины линии и ожидаемого
сдвига. При комнатной температуре и широком распределении размеров
эти признаки могут быть размыты; тогда более устойчивыми кандидатами
для диагностики являются поляризационный контраст и зависимость
от флюенса. Наличие доступного стоксова сдвига определяет возможность
спектрально отделить ФЛ от резонансного возбуждения. Если она отсутствует,
следует использовать дифференциальное пропускание, временное разделение
сигналов или резонансную спектроскопию с подавлением лазерного фона.

\paragraph{Нелинейное возбуждение и релаксация.}
При наличии фемтосекундного источника следует измерять зависимость
КТ-селективного сигнала от флюенса для обеих поляризаций, сохраняя
длительность, спектр и площадь пятна. Наиболее прямое дополнение --
спектроскопия «накачка--зондирование» с регистрацией изменения пропускания
в области экситонного перехода и его временной кинетики. В экспериментах
с КТ--Au такой метод позволяет выявлять быстрый перенос заряда,
отсутствующий в рассматриваемой электромагнитной модели~\cite{okuhata2020}.
Величину просветления нельзя без калибровки считать населённостью:
необходимо отделить экситонный вклад от плазмонного фона, изменения формы
линий и возможных многоэкситонных процессов. Собственная динамика золота
также зависит от флюенса и длины волны накачки~\cite{yu2011}, поэтому
контроль с одними МНЧ обязателен. Расчёт одного возбуждающего импульса
не является готовым расчётом сигнала «накачка--зондирование».

Более доступное дополнение к стационарной ФЛ -- измерение её времени
жизни методом времякоррелированного счёта одиночных фотонов (TCSPC).
Его временного разрешения достаточно лишь для компонент, медленных
относительно аппаратной функции; фемтосекундное тушение требует
быстрого зондирования. Для отдельной компоненты распада
$\tau^{-1}=\gamma_{\rm r}+\gamma_{\rm nr}$ и
$\Phi_{\rm PL}=\gamma_{\rm r}/(\gamma_{\rm r}+\gamma_{\rm nr})$,
поэтому сокращение $\tau$ само по себе не доказывает усиления излучения.
Нужны также квантовый выход или независимая калибровка излучательной
и безызлучательной долей; значимость расстояния для тушения показана
в~\cite{pons2007}. После окончания когерентной стадии и при независимых
центрах число зарегистрированных фотонов имеет вид
\begin{equation}
  N_{\rm det}\propto
  \sum_i \eta_{{\rm coll},i}\,\Phi_{{\rm PL},i}\,
  P_{{\rm exc},i}(\mathcal F),
  \label{eq:pl-readout}
\end{equation}
где $\eta_{{\rm coll},i}$ учитывает сбор и регистрацию.
Следовательно, порог $\rho_{ee}=0.5$ нельзя заменять точкой половинной
яркости ФЛ без проверки модели считывания. Для сопоставления с
$t_{\rm read}=\SI{3.6}{\pico\second}$ нужно восстановить населённость
именно на этой задержке, а не интеграл излучения за всё время.

\paragraph{Определение преобладающей доли.}
Для разреженного ансамбля независимо отвечающих комплексов совместную
обработку можно строить в виде
\begin{equation}
  Y_j=b_j+N_{\rm pair}\sum_\nu f_\nu
  \int q_{\nu j}(\boldsymbol\zeta)\,
              p_\nu(\boldsymbol\zeta)\,d\boldsymbol\zeta,
  \qquad f_\nu\ge0,\quad \sum_\nu f_\nu=1.
  \label{eq:ensemble-inference}
\end{equation}
Здесь $j$ обозначает энергию, поляризацию, флюенс и задержку измерения,
$b_j$ -- фон, $f_\nu$ -- численная доля геометрии $\nu$,
$p_\nu$ -- нормированное распределение зазоров, размеров, ориентаций
и параметров перехода, а $q_{\nu j}$ -- сигнал одной пары с учётом
конкретного способа регистрации. Свободные КТ и непарные МНЧ включаются
как отдельные компоненты с собственными концентрациями.
В частности, пять каналов табл.~\ref{tab:channels} не следует считать
пятью независимо подгоняемыми геометрическими долями: торцевая КТ
при двух поляризациях принадлежит одной и той же подгруппе.

Одни и те же $f_\nu$ должны описывать все поляризационные, спектральные
и флюенсные серии. О преобладании торцевых комплексов по числу можно
говорить, если после учёта неопределённостей нижняя граница оценки
$f_{\rm tip}$ превышает $0.5$; оптическое доминирование этого не гарантирует.
Небольшая доля особенно ярких пар может определять почти всю ФЛ.
Без калибровки отклика на одну пару, квантового выхода и распределений
параметров корректнее сообщать доли оптического сигнала либо диапазон
совместимых составов. Практическая проверка различимости -- совместная
подгонка эталонных смесей с известными соотношениями компонентов и анализ
того, допускают ли ошибки измерения альтернативный состав.

Изотропный ансамбль утрачивает информацию о выделенном направлении
в поляризационных измерениях. Для закреплённого слоя одного вращения
поляризации в его плоскости недостаточно, чтобы различить все положения
КТ над и под стержнем; требуются дополнительные проекции поля, например
при наклонном падении с $s$- и $p$-поляризациями, и учёт подложки.
Два эквивалентных конца симметричного сфероида в рассматриваемой модели
неразличимы. Поэтому цель предложенной схемы -- оценка преобладания
классов «у вершины» и «сбоку» при контролируемой ориентации и заданной
модели ансамбля, а не однозначное восстановление координат каждой КТ.

\subsection{Получение ориентированных комплексов доступными методами}
\label{sec:experiment-assembly}

Получение требуемого образца включает две операции: преимущественное
присоединение КТ к заданной области МНЧ и ориентацию уже сформированных
комплексов. Одинаковая ориентация стержней не обеспечивает автоматически
торцевого расположения КТ; простое смешивание двух коллоидов не обеспечивает
ни того, ни другого.

Прямой пример химического управления положением КТ на золотом наностержне
дан Nepal и соавторами~\cite{nepal2013}. Их подход основан
на функционализации стержней аминоалкантиолами и связывании
с функционализированными КТ через тиольную химию. В дополнительных
материалах описаны режимы преимущественной функционализации концов
и всей поверхности: для 11-амино-1-ундекантиола меняются концентрация
CTAB и длительность обработки. Используются водные растворы,
перемешивание и выдержка при комнатной температуре. Это подходящая
основа для серии «преимущественно торцевое присоединение -- распределённое
присоединение». Второй образец не следует называть чисто боковым без
структурной проверки. Уменьшение концентрации реагирующих частиц
и подбор их соотношения позволяют ограничивать число КТ на одном
стержне; сохранение коллоидной устойчивости проверяется по спектрам
и контролю агрегации. Длину линкера нельзя автоматически приравнивать
к $g$: в расстояние входят оболочки обеих частиц и конформация связки.

Другой доступный химический принцип -- сборка посредством
функционализированных ДНК и биоспецифических связей. Cohen-Hoshen
и соавторы получили комплексы CdSe/ZnS--Au с контролем состава
и структуры и исследовали их поляризационный отклик~\cite{cohenhoshen2012}.
Это пример организации полупроводниковых и металлических частиц,
но не доказательство селективного присоединения КТ к вершине сфероида.
Для наностержня такой маршрут нужно сочетать с селективной
функционализацией его поверхности. Биомолекулярный разделитель обычно
задаёт другую геометрию и среду, чем короткая органическая связка,
поэтому для конкретного комплекса требуется повторный расчёт.

Для макроскопической ориентации наиболее простой исходный вариант --
введение комплексов в прозрачную полимерную плёнку с последующим
одноосным растяжением. Ориентация золотых стержней таким способом
показана в ПВС~\cite{vanderzande1999}, в том числе для частиц на его
поверхности~\cite{murphy2005}. Объединение этого приёма с торцевой
функционализацией~\cite{nepal2013} предлагается здесь как маршрут
для проверки, а не как уже продемонстрированная в этих работах
сборка рассматриваемой пары с зазором \SI{1}{\nano\meter}.
После растяжения необходимо подтвердить сохранность связок,
распределения зазоров и спектров КТ. Разреженная плёнка предпочтительна
для сравнения с моделью одиночной пары, поскольку уменьшает влияние
взаимодействия соседних МНЧ.

Пример направленной самосборки именно массива МНЧ с КТ представлен
в~\cite{peng2014}: CTAB-стабилизированные золотые стержни образуют
вертикально ориентированный монослой, над которым размещаются
CdSe/ZnS КТ. Однако полный протокол включает напыление разделяющего
слоя SiO$_2$, а оптический отклик относится к плотному массиву.
Этот способ полезен как подтверждение возможности получения
ориентированной структуры, но требует соответствующего оборудования
и не воспроизводит автоматически ни разреженную бинарную систему,
ни её оптимальный зазор. Для обычной химической лаборатории
первым выбором остаются сборка в растворе и ориентация в полимере.

Постоянное электрическое поле может ориентировать анизотропную частицу
за счёт индуцированного диполя. В простейшем случае энергия ориентации
содержит $-\Delta\alpha(0)E_{\rm dc}^2\cos^2\theta/2$; заметный порядок
возникает при конкуренции этой энергии с тепловым вращением.
Но это ориентирует ось МНЧ и не выбирает место химического связывания КТ.
Непосредственно продемонстрированный для коротких золотых стержней
вариант -- внешнее модулированное электрическое поле с контролем
ориентации по поляризованной экстинкции~\cite{zijlstra2012}.
Для проводящих коллоидных сред постоянное поле дополнительно связано
с переносом частиц и электродными процессами; его эффективность
для конкретных комплексов нужно устанавливать отдельно. Практически
разумно сначала сформировать химически связанные пары, затем
ориентировать их полем и закрепить в матрице. Требуемые напряжённость,
частота и время фиксации из оптической модели данной работы не следуют.

Лазерный пинцет можно рассматривать как дополнительную возможность
для лаборатории, уже располагающей такой установкой. Pelton
и соавторы~\cite{pelton2006} продемонстрировали захват и ориентацию
отдельных золотых наностержней при выборе длины волны с длинноволновой
стороны продольного плазмонного резонанса. Это доказывает возможность
управления металлической частицей, но не адресного закрепления
четырёхнанометровой КТ у её конца и не массовой сборки массива.
С учётом локального нагрева и ограниченной производительности
пинцет не является основным предлагаемым способом изготовления.

Для предпочтительного торцевого канала целевой образец представляет
собой разреженный слой стержней с общим направлением осей и
преимущественным присоединением КТ к концам. Если оси лежат в плоскости
слоя, нормальное падение света позволяет ориентировать электрическое
поле вдоль них и затем повернуть его на $90^\circ$ для контроля.
Это удобная экспериментальная реализация заданного в расчёте поля,
а не найденная моделью зависимость от направления распространения.
Для вертикальных стержней продольное поле, напротив, требует другой
геометрии освещения. Сопоставление измеренных спектров, нелинейных
кривых и кинетики для торцевого и распределённого присоединения
позволит проверить, сохраняется ли рассчитанное преимущество после
учёта структуры и потерь конкретного образца.
